%% groth16.tex -- Construction of the Groth16 SNARK
%% %% groth16.tex -- Construction of the Groth16 SNARK
%% %% groth16.tex -- Construction of the Groth16 SNARK
%% %% groth16.tex -- Construction of the Groth16 SNARK
%% \input{groth16} inside zero_knowledge.tex

%% ============================================================
\section{Groth16: A Pairing-Based SNARK}
%% ============================================================

Groth16 is a preprocessing SNARK for arithmetic circuit satisfiability. Its proof is only three group elements, verification is a single pairing equation, and it achieves perfect completeness and perfect zero knowledge. Soundness holds against adversaries that interact with the system via generic group operations.

We build the system in three layers: encode computation as a polynomial divisibility problem (Quadratic Arithmetic Program, QAP), construct an information-theoretic proof for it (Non-Interactive Linear Proof, NILP), then compile the NILP into group elements using pairings.

%% -----------------------------------------------------------
\subsection{Notation for Pairing Groups}
%% -----------------------------------------------------------

Let $(p, \G_1, \G_2, \GT, e, g, h)$ be an asymmetric bilinear group of prime order $p$. The generator of $\G_1$ is $g$, the generator of $\G_2$ is $h$, and $e : \G_1 \times \G_2 \to \GT$ is a bilinear pairing satisfying $e(g^a, h^b) = e(g,h)^{ab}$.

We represent group elements by their discrete logarithms using the bracket notation
\[
  \gone{a} \deq g^a \in \G_1, \qquad
  \gtwo{a} \deq h^a \in \G_2, \qquad
  \gT{a}   \deq e(g,h)^a \in \GT, \qquad a \in \Z_p.
\]
This notation makes linear algebra readable. For a known scalar $c$ or known matrix $M$:
\[
  \gone{a} + \gone{b} = \gone{a+b}, \qquad
  c \cdot \gone{a} = \gone{ca}, \qquad
  M \cdot \gone{\mathbf{a}} = \gone{M\mathbf{a}}.
\]
The pairing provides one level of multiplication across the two source groups:
\[
  \gone{a} \cdot \gtwo{b} \deq e\!\left(g^a, h^b\right) = \gT{ab}.
\]


%% -----------------------------------------------------------
\subsection{From Arithmetic Circuits to Quadratic Arithmetic Programs}
%% -----------------------------------------------------------

\subsubsection*{Arithmetic circuits and R1CS}

An arithmetic circuit over a finite field $\F_p$ consists of addition and multiplication gates over field elements. Any such circuit can be \emph{flattened}: every gate is rewritten as a single multiplicative constraint
\[
  \bigl(\textstyle\sum_i A_{q,i}\, a_i\bigr)
  \cdot
  \bigl(\textstyle\sum_i B_{q,i}\, a_i\bigr)
  =
  \textstyle\sum_i C_{q,i}\, a_i
  \qquad (q = 1,\ldots, n),
\]
where $z = (a_0, a_1, \ldots, a_m) \in \F_p^{m+1}$ is the \emph{assignment vector}: $a_0 = 1$ is a constant, $\phi = (a_1, \ldots, a_\ell)$ is the public statement, and $w = (a_{\ell+1}, \ldots, a_m)$ is the secret witness. The $n$ constraints are collected into three matrices $A, B, C \in \F_p^{n \times (m+1)}$ --- this is the \emph{rank-1 constraint system} (R1CS). Arithmetic circuit satisfiability is NP-complete: any NP statement can be compiled into an R1CS instance over a suitable field, so a proof system for R1CS is a proof system for all of NP.

\subsubsection*{Quadratic arithmetic programs}

Checking $n$ separate scalar equations one by one would be costly. The key algebraic trick is to compress them into a single polynomial divisibility test.

Fix $n$ distinct field elements $r_1, \ldots, r_n \in \F_p$, one for each constraint. Define the vanishing polynomial
\[
  t(X) = \prod_{q=1}^{n}(X - r_q).
\]
For each variable index $i \in \{0, \ldots, m\}$, interpolate three polynomials $u_i, v_i, w_i \in \F_p[X]$ of degree at most $n-1$ such that
\[
  u_i(r_q) = A_{q,i}, \qquad v_i(r_q) = B_{q,i}, \qquad w_i(r_q) = C_{q,i} \qquad \forall\, q.
\]
\todo{Add a visualisation showing how the columns of the $A$, $B$, $C$ matrices are read off and interpolated into the polynomials $u_i$, $v_i$, $w_i$.}
Given an assignment $z$, define the aggregated polynomials
\[
  U(X) = \sum_{i=0}^{m} a_i\, u_i(X), \qquad
  V(X) = \sum_{i=0}^{m} a_i\, v_i(X), \qquad
  W(X) = \sum_{i=0}^{m} a_i\, w_i(X).
\]
At each constraint point $r_q$, the R1CS equation $U(r_q)\cdot V(r_q) = W(r_q)$ must hold. Since $t(X)$ vanishes precisely at $r_1, \ldots, r_n$, this is equivalent to
\[
  \boxed{t(X) \;\Big|\; U(X)\cdot V(X) - W(X).}
\]
Equivalently, there exists a degree-$(n-2)$ quotient polynomial $h(X)$ such that
\[
  U(X)\cdot V(X) = W(X) + h(X)\,t(X).
\]
This is the \emph{quadratic arithmetic program} (QAP) relation for the circuit. The advantage over R1CS is that $n$ separate equations have been replaced by one polynomial identity checkable at a single random point --- exactly the Schwartz--Zippel trick.

%% -----------------------------------------------------------
\subsection{A Non-Interactive Linear Proof for QAP}
%% -----------------------------------------------------------

A \emph{non-interactive linear proof} (NILP) is an information-theoretic proof system where the prover's messages are linear functions of a common reference string $\sigma$. We now construct such a system for the QAP relation.

\todo{Explain why this NILP is non-interactive despite stemming from a Linear Interactive Proof (LIP): the CRS plays the role of the verifier's challenges, collapsing the interaction into a single message. Discuss the connection between 2-move LIPs and NILPs more carefully.}

\subsubsection*{Setup}

Pick a random trapdoor $\tau = (\alpha, \beta, \gamma, \delta, x) \in \F_p^{*5}$ and define the common reference string $\sigma$ as the collection of field evaluations
\[
\sigma =
\left(
\begin{aligned}
&\alpha,\; \beta,\; \gamma,\; \delta,\;\{x^i\}_{i=0}^{n-1},\\[2pt]
&\left\{\frac{\beta u_i(x)+\alpha v_i(x)+w_i(x)}{\gamma}\right\}_{i=0}^{\ell},\\[4pt]
&\left\{\frac{\beta u_i(x)+\alpha v_i(x)+w_i(x)}{\delta}\right\}_{i=\ell+1}^{m},\\[4pt]
&\left\{\frac{x^i\, t(x)}{\delta}\right\}_{i=0}^{n-2}
\end{aligned}
\right).
\]
The first block enables evaluation of arbitrary polynomials at $x$. The second block encodes the public witness terms, scaled by $\gamma$. The third encodes the private witness terms, scaled by $\delta$. The last encodes powers of $x$ weighted by $t(x)/\delta$, allowing the prover to compute $h(x)\,t(x)/\delta$ as a linear combination.

\todo{Discuss what $\alpha$, $\beta$, $\gamma$, and $\delta$ each do exactly: why they appear in the CRS, how they enforce consistency between $A$, $B$, $C$, and why the $\gamma$/$\delta$ separation is necessary.}

\subsubsection*{Prover}

Given a satisfying assignment $(a_1, \ldots, a_m)$, the prover computes $h(X)$ from the QAP relation, then picks two random blinding scalars $r, s \leftarrow \F_p$ and sets
\begin{align*}
  A &\deq \alpha + U(x) + r\delta, \\
  B &\deq \beta  + V(x) + s\delta, \\
  C &\deq \frac{\displaystyle\sum_{i=\ell+1}^{m} a_i\bigl(\beta u_i(x)+\alpha v_i(x)+w_i(x)\bigr) + h(x)\,t(x)}{\delta} + As + rB - rs\delta.
\end{align*}
The proof is $\pi = (A, B, C)$.

\subsubsection*{Verifier}

The verifier checks the single equation
\[
  A \cdot B
  \;\stackrel{?}{=}\;
  \underbrace{\alpha\beta}_{\text{(i)}}
  \;+\;
  \underbrace{\sum_{i=0}^{\ell} a_i \cdot \frac{\beta u_i(x)+\alpha v_i(x)+w_i(x)}{\gamma}}_{\text{(ii)}} \cdot\, \gamma
  \;+\;
  \underbrace{C \cdot \delta}_{\text{(iii)}}.
\]
The three terms on the right play distinct roles.
\begin{itemize}[noitemsep, leftmargin=2em]
  \item[(i)] $\alpha\beta$ is a fixed anchor from the setup. Its presence forces $A$ and $B$ to contain genuine $\alpha$ and $\beta$ offsets respectively, preventing a prover from constructing $A$ and $B$ independently of the structure of $\sigma$.
  \item[(ii)] The sum over $i = 0, \ldots, \ell$ incorporates the \emph{public inputs} $a_0, \ldots, a_\ell$ --- the statement $\phi$ that the verifier knows. These are precomputed in the CRS scaled by $\gamma$ and $\gamma$ cancels out, leaving the verifier to evaluate this term using only public data. Note that $\alpha$ and $\beta$ also appear here, but they are \emph{setup parameters} embedded in the precomputed CRS entries, not the circuit's inputs.
  \item[(iii)] $C \cdot \delta$ is where the private witness $a_{\ell+1}, \ldots, a_m$ is encoded. The $\delta$ scaling keeps it algebraically separated from term (ii), preventing the prover from mixing public and private contributions.
\end{itemize}

\subsubsection*{Simulator}

Given the trapdoor $\tau$, the simulator ignores the witness entirely. It picks $A, B \leftarrow \F_p$ uniformly at random and computes
\[
  C \deq \frac{A B \;-\; \alpha\beta \;-\; \displaystyle\sum_{i=0}^{\ell} a_i\bigl(\beta u_i(x)+\alpha v_i(x)+w_i(x)\bigr)}{\delta}.
\]

\subsubsection*{Why this achieves zero knowledge}

The blinding scalars $r$ and $s$ make $A$ and $B$ uniformly distributed in $\F_p$, independent of the witness. Given uniform $A$ and $B$, the verification equation uniquely determines $C$. The real distribution of $(A, B, C)$ is therefore identical to the simulated one: two uniform field elements followed by their forced companion. This is perfect zero knowledge.

\subsubsection*{Why this is sound against affine strategies}

An affine prover strategy computes $(A, B, C)$ as linear combinations of $\sigma$. Viewing the verification equation as a polynomial identity in the formal indeterminates $\alpha, \beta, \gamma, \delta, x$ and applying Schwartz--Zippel, any affine strategy that passes verification with non-negligible probability must satisfy the identity formally. A careful analysis of the monomials (particularly those involving $\alpha\beta$, $1/\gamma^2$, and $1/\delta^2$) forces the coefficients to encode a consistent assignment, from which the witness $(a_{\ell+1}, \ldots, a_m)$ can be extracted.

\todo{Add the full monomial analysis proof here, following Theorem~1 of the Groth (2016) paper.}

%% -----------------------------------------------------------
\subsection{Compiling to Pairings: The Generic Group Model}
%% -----------------------------------------------------------

The NILP works over a field, but we need a cryptographic argument --- the prover should not be able to read off $x$ from $\sigma$. The compilation step replaces every field element in $\sigma$ with its bracket encoding, and $(A, B, C)$ become $(\gone{A}, \gtwo{B}, \gone{C})$.

The compiled verification equation is:
\[
  \gone{A} \cdot \gtwo{B}
  \;=\;
  \gone{\alpha}\cdot\gtwo{\beta}
  \;+\;
  \sum_{i=0}^{\ell} a_i\,
    \left[\frac{\beta u_i(x)+\alpha v_i(x)+w_i(x)}{\gamma}\right]_{\!1}
    \cdot\, \gtwo{\gamma}
  \;+\;
  \gone{C}\cdot\gtwo{\delta}.
\]
The prover computes $\gone{A}$, $\gone{C}$, $\gtwo{B}$ as linear combinations of the CRS elements --- this is exactly the NILP prover running in the exponent. The verifier checks one pairing product equation using three pairings in total.

\textbf{Lifting soundness via the generic group model.} In the generic group model, a prover can only use the group elements in $\sigma$ through addition (and scalar multiplication) --- it cannot read their discrete logarithms. Any computation it performs on $\gone{\sigma_1}$ and $\gtwo{\sigma_2}$ is therefore a sequence of affine operations, corresponding exactly to choosing a proof matrix $\Pi$ and computing $\Pi\sigma$. This is precisely an affine strategy in the NILP sense. Combined with a \emph{disclosure-freeness} property of the CRS --- which guarantees that the structure of $\sigma$ cannot be exploited to choose a non-affine strategy --- the NILP's soundness against affine strategies lifts to soundness against all generic adversaries.

%% -----------------------------------------------------------
\subsection*{Summary: The Groth16 Proof System}
%% -----------------------------------------------------------

\begin{tcolorbox}[enhanced, colback=blue!3!white, colframe=blue!45!black,
                  arc=2.5pt, boxrule=1pt, left=6pt, right=6pt, top=5pt, bottom=5pt]
\small
\textbf{Setup.} Sample $\tau = (\alpha,\beta,\gamma,\delta,x) \leftarrow \F_p^{*5}$.
Publish $\sigma = \bigl([\sigma_1]_1,\, [\sigma_2]_2\bigr)$ as described above. Keep $\tau$ secret.

\medskip
\textbf{Prove.} Given $(a_1,\ldots,a_m)$ satisfying the circuit, pick $r,s \leftarrow \F_p$ and output
$\pi = \bigl(\gone{A},\, \gone{C},\, \gtwo{B}\bigr) \in \G_1^2 \times \G_2$.

\medskip
\textbf{Verify.} Check:
$\gone{A} \cdot \gtwo{B} = \gone{\alpha}\cdot\gtwo{\beta} + \displaystyle\sum_{i=0}^{\ell} a_i\,\gone{(\cdot)_i}\cdot\gtwo{\gamma} + \gone{C}\cdot\gtwo{\delta}$.

\medskip
\textbf{Proof size:} \quad $2$ elements in $\G_1$, \quad $1$ element in $\G_2$ \quad $= 3$ group elements.\\
\textbf{Verification:} \quad $1$ pairing product equation requiring $3$ online pairings:
$e(\gone{A}, \gtwo{B})$,\; $e\!\left(\textstyle\sum a_i K_i,\, \gtwo{\gamma}\right)$,\; $e(\gone{C}, \gtwo{\delta})$.
The fourth term $\gT{\alpha\beta}$ is precomputed and stored in the verifier key, so no extra pairing is needed for it.
\end{tcolorbox}
 inside zero_knowledge.tex

%% ============================================================
\section{Groth16: A Pairing-Based SNARK}
%% ============================================================

Groth16 is a preprocessing SNARK for arithmetic circuit satisfiability. Its proof is only three group elements, verification is a single pairing equation, and it achieves perfect completeness and perfect zero knowledge. Soundness holds against adversaries that interact with the system via generic group operations.

We build the system in three layers: encode computation as a polynomial divisibility problem (Quadratic Arithmetic Program, QAP), construct an information-theoretic proof for it (Non-Interactive Linear Proof, NILP), then compile the NILP into group elements using pairings.

%% -----------------------------------------------------------
\subsection{Notation for Pairing Groups}
%% -----------------------------------------------------------

Let $(p, \G_1, \G_2, \GT, e, g, h)$ be an asymmetric bilinear group of prime order $p$. The generator of $\G_1$ is $g$, the generator of $\G_2$ is $h$, and $e : \G_1 \times \G_2 \to \GT$ is a bilinear pairing satisfying $e(g^a, h^b) = e(g,h)^{ab}$.

We represent group elements by their discrete logarithms using the bracket notation
\[
  \gone{a} \deq g^a \in \G_1, \qquad
  \gtwo{a} \deq h^a \in \G_2, \qquad
  \gT{a}   \deq e(g,h)^a \in \GT, \qquad a \in \Z_p.
\]
This notation makes linear algebra readable. For a known scalar $c$ or known matrix $M$:
\[
  \gone{a} + \gone{b} = \gone{a+b}, \qquad
  c \cdot \gone{a} = \gone{ca}, \qquad
  M \cdot \gone{\mathbf{a}} = \gone{M\mathbf{a}}.
\]
The pairing provides one level of multiplication across the two source groups:
\[
  \gone{a} \cdot \gtwo{b} \deq e\!\left(g^a, h^b\right) = \gT{ab}.
\]


%% -----------------------------------------------------------
\subsection{From Arithmetic Circuits to Quadratic Arithmetic Programs}
%% -----------------------------------------------------------

\subsubsection*{Arithmetic circuits and R1CS}

An arithmetic circuit over a finite field $\F_p$ consists of addition and multiplication gates over field elements. Any such circuit can be \emph{flattened}: every gate is rewritten as a single multiplicative constraint
\[
  \bigl(\textstyle\sum_i A_{q,i}\, a_i\bigr)
  \cdot
  \bigl(\textstyle\sum_i B_{q,i}\, a_i\bigr)
  =
  \textstyle\sum_i C_{q,i}\, a_i
  \qquad (q = 1,\ldots, n),
\]
where $z = (a_0, a_1, \ldots, a_m) \in \F_p^{m+1}$ is the \emph{assignment vector}: $a_0 = 1$ is a constant, $\phi = (a_1, \ldots, a_\ell)$ is the public statement, and $w = (a_{\ell+1}, \ldots, a_m)$ is the secret witness. The $n$ constraints are collected into three matrices $A, B, C \in \F_p^{n \times (m+1)}$ --- this is the \emph{rank-1 constraint system} (R1CS). Arithmetic circuit satisfiability is NP-complete: any NP statement can be compiled into an R1CS instance over a suitable field, so a proof system for R1CS is a proof system for all of NP.

\subsubsection*{Quadratic arithmetic programs}

Checking $n$ separate scalar equations one by one would be costly. The key algebraic trick is to compress them into a single polynomial divisibility test.

Fix $n$ distinct field elements $r_1, \ldots, r_n \in \F_p$, one for each constraint. Define the vanishing polynomial
\[
  t(X) = \prod_{q=1}^{n}(X - r_q).
\]
For each variable index $i \in \{0, \ldots, m\}$, interpolate three polynomials $u_i, v_i, w_i \in \F_p[X]$ of degree at most $n-1$ such that
\[
  u_i(r_q) = A_{q,i}, \qquad v_i(r_q) = B_{q,i}, \qquad w_i(r_q) = C_{q,i} \qquad \forall\, q.
\]
\todo{Add a visualisation showing how the columns of the $A$, $B$, $C$ matrices are read off and interpolated into the polynomials $u_i$, $v_i$, $w_i$.}
Given an assignment $z$, define the aggregated polynomials
\[
  U(X) = \sum_{i=0}^{m} a_i\, u_i(X), \qquad
  V(X) = \sum_{i=0}^{m} a_i\, v_i(X), \qquad
  W(X) = \sum_{i=0}^{m} a_i\, w_i(X).
\]
At each constraint point $r_q$, the R1CS equation $U(r_q)\cdot V(r_q) = W(r_q)$ must hold. Since $t(X)$ vanishes precisely at $r_1, \ldots, r_n$, this is equivalent to
\[
  \boxed{t(X) \;\Big|\; U(X)\cdot V(X) - W(X).}
\]
Equivalently, there exists a degree-$(n-2)$ quotient polynomial $h(X)$ such that
\[
  U(X)\cdot V(X) = W(X) + h(X)\,t(X).
\]
This is the \emph{quadratic arithmetic program} (QAP) relation for the circuit. The advantage over R1CS is that $n$ separate equations have been replaced by one polynomial identity checkable at a single random point --- exactly the Schwartz--Zippel trick.

%% -----------------------------------------------------------
\subsection{A Non-Interactive Linear Proof for QAP}
%% -----------------------------------------------------------

A \emph{non-interactive linear proof} (NILP) is an information-theoretic proof system where the prover's messages are linear functions of a common reference string $\sigma$. We now construct such a system for the QAP relation.

\todo{Explain why this NILP is non-interactive despite stemming from a Linear Interactive Proof (LIP): the CRS plays the role of the verifier's challenges, collapsing the interaction into a single message. Discuss the connection between 2-move LIPs and NILPs more carefully.}

\subsubsection*{Setup}

Pick a random trapdoor $\tau = (\alpha, \beta, \gamma, \delta, x) \in \F_p^{*5}$ and define the common reference string $\sigma$ as the collection of field evaluations
\[
\sigma =
\left(
\begin{aligned}
&\alpha,\; \beta,\; \gamma,\; \delta,\;\{x^i\}_{i=0}^{n-1},\\[2pt]
&\left\{\frac{\beta u_i(x)+\alpha v_i(x)+w_i(x)}{\gamma}\right\}_{i=0}^{\ell},\\[4pt]
&\left\{\frac{\beta u_i(x)+\alpha v_i(x)+w_i(x)}{\delta}\right\}_{i=\ell+1}^{m},\\[4pt]
&\left\{\frac{x^i\, t(x)}{\delta}\right\}_{i=0}^{n-2}
\end{aligned}
\right).
\]
The first block enables evaluation of arbitrary polynomials at $x$. The second block encodes the public witness terms, scaled by $\gamma$. The third encodes the private witness terms, scaled by $\delta$. The last encodes powers of $x$ weighted by $t(x)/\delta$, allowing the prover to compute $h(x)\,t(x)/\delta$ as a linear combination.

\todo{Discuss what $\alpha$, $\beta$, $\gamma$, and $\delta$ each do exactly: why they appear in the CRS, how they enforce consistency between $A$, $B$, $C$, and why the $\gamma$/$\delta$ separation is necessary.}

\subsubsection*{Prover}

Given a satisfying assignment $(a_1, \ldots, a_m)$, the prover computes $h(X)$ from the QAP relation, then picks two random blinding scalars $r, s \leftarrow \F_p$ and sets
\begin{align*}
  A &\deq \alpha + U(x) + r\delta, \\
  B &\deq \beta  + V(x) + s\delta, \\
  C &\deq \frac{\displaystyle\sum_{i=\ell+1}^{m} a_i\bigl(\beta u_i(x)+\alpha v_i(x)+w_i(x)\bigr) + h(x)\,t(x)}{\delta} + As + rB - rs\delta.
\end{align*}
The proof is $\pi = (A, B, C)$.

\subsubsection*{Verifier}

The verifier checks the single equation
\[
  A \cdot B
  \;\stackrel{?}{=}\;
  \underbrace{\alpha\beta}_{\text{(i)}}
  \;+\;
  \underbrace{\sum_{i=0}^{\ell} a_i \cdot \frac{\beta u_i(x)+\alpha v_i(x)+w_i(x)}{\gamma}}_{\text{(ii)}} \cdot\, \gamma
  \;+\;
  \underbrace{C \cdot \delta}_{\text{(iii)}}.
\]
The three terms on the right play distinct roles.
\begin{itemize}[noitemsep, leftmargin=2em]
  \item[(i)] $\alpha\beta$ is a fixed anchor from the setup. Its presence forces $A$ and $B$ to contain genuine $\alpha$ and $\beta$ offsets respectively, preventing a prover from constructing $A$ and $B$ independently of the structure of $\sigma$.
  \item[(ii)] The sum over $i = 0, \ldots, \ell$ incorporates the \emph{public inputs} $a_0, \ldots, a_\ell$ --- the statement $\phi$ that the verifier knows. These are precomputed in the CRS scaled by $\gamma$ and $\gamma$ cancels out, leaving the verifier to evaluate this term using only public data. Note that $\alpha$ and $\beta$ also appear here, but they are \emph{setup parameters} embedded in the precomputed CRS entries, not the circuit's inputs.
  \item[(iii)] $C \cdot \delta$ is where the private witness $a_{\ell+1}, \ldots, a_m$ is encoded. The $\delta$ scaling keeps it algebraically separated from term (ii), preventing the prover from mixing public and private contributions.
\end{itemize}

\subsubsection*{Simulator}

Given the trapdoor $\tau$, the simulator ignores the witness entirely. It picks $A, B \leftarrow \F_p$ uniformly at random and computes
\[
  C \deq \frac{A B \;-\; \alpha\beta \;-\; \displaystyle\sum_{i=0}^{\ell} a_i\bigl(\beta u_i(x)+\alpha v_i(x)+w_i(x)\bigr)}{\delta}.
\]

\subsubsection*{Why this achieves zero knowledge}

The blinding scalars $r$ and $s$ make $A$ and $B$ uniformly distributed in $\F_p$, independent of the witness. Given uniform $A$ and $B$, the verification equation uniquely determines $C$. The real distribution of $(A, B, C)$ is therefore identical to the simulated one: two uniform field elements followed by their forced companion. This is perfect zero knowledge.

\subsubsection*{Why this is sound against affine strategies}

An affine prover strategy computes $(A, B, C)$ as linear combinations of $\sigma$. Viewing the verification equation as a polynomial identity in the formal indeterminates $\alpha, \beta, \gamma, \delta, x$ and applying Schwartz--Zippel, any affine strategy that passes verification with non-negligible probability must satisfy the identity formally. A careful analysis of the monomials (particularly those involving $\alpha\beta$, $1/\gamma^2$, and $1/\delta^2$) forces the coefficients to encode a consistent assignment, from which the witness $(a_{\ell+1}, \ldots, a_m)$ can be extracted.

\todo{Add the full monomial analysis proof here, following Theorem~1 of the Groth (2016) paper.}

%% -----------------------------------------------------------
\subsection{Compiling to Pairings: The Generic Group Model}
%% -----------------------------------------------------------

The NILP works over a field, but we need a cryptographic argument --- the prover should not be able to read off $x$ from $\sigma$. The compilation step replaces every field element in $\sigma$ with its bracket encoding, and $(A, B, C)$ become $(\gone{A}, \gtwo{B}, \gone{C})$.

The compiled verification equation is:
\[
  \gone{A} \cdot \gtwo{B}
  \;=\;
  \gone{\alpha}\cdot\gtwo{\beta}
  \;+\;
  \sum_{i=0}^{\ell} a_i\,
    \left[\frac{\beta u_i(x)+\alpha v_i(x)+w_i(x)}{\gamma}\right]_{\!1}
    \cdot\, \gtwo{\gamma}
  \;+\;
  \gone{C}\cdot\gtwo{\delta}.
\]
The prover computes $\gone{A}$, $\gone{C}$, $\gtwo{B}$ as linear combinations of the CRS elements --- this is exactly the NILP prover running in the exponent. The verifier checks one pairing product equation using three pairings in total.

\textbf{Lifting soundness via the generic group model.} In the generic group model, a prover can only use the group elements in $\sigma$ through addition (and scalar multiplication) --- it cannot read their discrete logarithms. Any computation it performs on $\gone{\sigma_1}$ and $\gtwo{\sigma_2}$ is therefore a sequence of affine operations, corresponding exactly to choosing a proof matrix $\Pi$ and computing $\Pi\sigma$. This is precisely an affine strategy in the NILP sense. Combined with a \emph{disclosure-freeness} property of the CRS --- which guarantees that the structure of $\sigma$ cannot be exploited to choose a non-affine strategy --- the NILP's soundness against affine strategies lifts to soundness against all generic adversaries.

%% -----------------------------------------------------------
\subsection*{Summary: The Groth16 Proof System}
%% -----------------------------------------------------------

\begin{tcolorbox}[enhanced, colback=blue!3!white, colframe=blue!45!black,
                  arc=2.5pt, boxrule=1pt, left=6pt, right=6pt, top=5pt, bottom=5pt]
\small
\textbf{Setup.} Sample $\tau = (\alpha,\beta,\gamma,\delta,x) \leftarrow \F_p^{*5}$.
Publish $\sigma = \bigl([\sigma_1]_1,\, [\sigma_2]_2\bigr)$ as described above. Keep $\tau$ secret.

\medskip
\textbf{Prove.} Given $(a_1,\ldots,a_m)$ satisfying the circuit, pick $r,s \leftarrow \F_p$ and output
$\pi = \bigl(\gone{A},\, \gone{C},\, \gtwo{B}\bigr) \in \G_1^2 \times \G_2$.

\medskip
\textbf{Verify.} Check:
$\gone{A} \cdot \gtwo{B} = \gone{\alpha}\cdot\gtwo{\beta} + \displaystyle\sum_{i=0}^{\ell} a_i\,\gone{(\cdot)_i}\cdot\gtwo{\gamma} + \gone{C}\cdot\gtwo{\delta}$.

\medskip
\textbf{Proof size:} \quad $2$ elements in $\G_1$, \quad $1$ element in $\G_2$ \quad $= 3$ group elements.\\
\textbf{Verification:} \quad $1$ pairing product equation requiring $3$ online pairings:
$e(\gone{A}, \gtwo{B})$,\; $e\!\left(\textstyle\sum a_i K_i,\, \gtwo{\gamma}\right)$,\; $e(\gone{C}, \gtwo{\delta})$.
The fourth term $\gT{\alpha\beta}$ is precomputed and stored in the verifier key, so no extra pairing is needed for it.
\end{tcolorbox}
 inside zero_knowledge.tex

%% ============================================================
\section{Groth16: A Pairing-Based SNARK}
%% ============================================================

Groth16 is a preprocessing SNARK for arithmetic circuit satisfiability. Its proof is only three group elements, verification is a single pairing equation, and it achieves perfect completeness and perfect zero knowledge. Soundness holds against adversaries that interact with the system via generic group operations.

We build the system in three layers: encode computation as a polynomial divisibility problem (Quadratic Arithmetic Program, QAP), construct an information-theoretic proof for it (Non-Interactive Linear Proof, NILP), then compile the NILP into group elements using pairings.

%% -----------------------------------------------------------
\subsection{Notation for Pairing Groups}
%% -----------------------------------------------------------

Let $(p, \G_1, \G_2, \GT, e, g, h)$ be an asymmetric bilinear group of prime order $p$. The generator of $\G_1$ is $g$, the generator of $\G_2$ is $h$, and $e : \G_1 \times \G_2 \to \GT$ is a bilinear pairing satisfying $e(g^a, h^b) = e(g,h)^{ab}$.

We represent group elements by their discrete logarithms using the bracket notation
\[
  \gone{a} \deq g^a \in \G_1, \qquad
  \gtwo{a} \deq h^a \in \G_2, \qquad
  \gT{a}   \deq e(g,h)^a \in \GT, \qquad a \in \Z_p.
\]
This notation makes linear algebra readable. For a known scalar $c$ or known matrix $M$:
\[
  \gone{a} + \gone{b} = \gone{a+b}, \qquad
  c \cdot \gone{a} = \gone{ca}, \qquad
  M \cdot \gone{\mathbf{a}} = \gone{M\mathbf{a}}.
\]
The pairing provides one level of multiplication across the two source groups:
\[
  \gone{a} \cdot \gtwo{b} \deq e\!\left(g^a, h^b\right) = \gT{ab}.
\]


%% -----------------------------------------------------------
\subsection{From Arithmetic Circuits to Quadratic Arithmetic Programs}
%% -----------------------------------------------------------

\subsubsection*{Arithmetic circuits and R1CS}

An arithmetic circuit over a finite field $\F_p$ consists of addition and multiplication gates over field elements. Any such circuit can be \emph{flattened}: every gate is rewritten as a single multiplicative constraint
\[
  \bigl(\textstyle\sum_i A_{q,i}\, a_i\bigr)
  \cdot
  \bigl(\textstyle\sum_i B_{q,i}\, a_i\bigr)
  =
  \textstyle\sum_i C_{q,i}\, a_i
  \qquad (q = 1,\ldots, n),
\]
where $z = (a_0, a_1, \ldots, a_m) \in \F_p^{m+1}$ is the \emph{assignment vector}: $a_0 = 1$ is a constant, $\phi = (a_1, \ldots, a_\ell)$ is the public statement, and $w = (a_{\ell+1}, \ldots, a_m)$ is the secret witness. The $n$ constraints are collected into three matrices $A, B, C \in \F_p^{n \times (m+1)}$ --- this is the \emph{rank-1 constraint system} (R1CS). Arithmetic circuit satisfiability is NP-complete: any NP statement can be compiled into an R1CS instance over a suitable field, so a proof system for R1CS is a proof system for all of NP.

\subsubsection*{Quadratic arithmetic programs}

Checking $n$ separate scalar equations one by one would be costly. The key algebraic trick is to compress them into a single polynomial divisibility test.

Fix $n$ distinct field elements $r_1, \ldots, r_n \in \F_p$, one for each constraint. Define the vanishing polynomial
\[
  t(X) = \prod_{q=1}^{n}(X - r_q).
\]
For each variable index $i \in \{0, \ldots, m\}$, interpolate three polynomials $u_i, v_i, w_i \in \F_p[X]$ of degree at most $n-1$ such that
\[
  u_i(r_q) = A_{q,i}, \qquad v_i(r_q) = B_{q,i}, \qquad w_i(r_q) = C_{q,i} \qquad \forall\, q.
\]
\todo{Add a visualisation showing how the columns of the $A$, $B$, $C$ matrices are read off and interpolated into the polynomials $u_i$, $v_i$, $w_i$.}
Given an assignment $z$, define the aggregated polynomials
\[
  U(X) = \sum_{i=0}^{m} a_i\, u_i(X), \qquad
  V(X) = \sum_{i=0}^{m} a_i\, v_i(X), \qquad
  W(X) = \sum_{i=0}^{m} a_i\, w_i(X).
\]
At each constraint point $r_q$, the R1CS equation $U(r_q)\cdot V(r_q) = W(r_q)$ must hold. Since $t(X)$ vanishes precisely at $r_1, \ldots, r_n$, this is equivalent to
\[
  \boxed{t(X) \;\Big|\; U(X)\cdot V(X) - W(X).}
\]
Equivalently, there exists a degree-$(n-2)$ quotient polynomial $h(X)$ such that
\[
  U(X)\cdot V(X) = W(X) + h(X)\,t(X).
\]
This is the \emph{quadratic arithmetic program} (QAP) relation for the circuit. The advantage over R1CS is that $n$ separate equations have been replaced by one polynomial identity checkable at a single random point --- exactly the Schwartz--Zippel trick.

%% -----------------------------------------------------------
\subsection{A Non-Interactive Linear Proof for QAP}
%% -----------------------------------------------------------

A \emph{non-interactive linear proof} (NILP) is an information-theoretic proof system where the prover's messages are linear functions of a common reference string $\sigma$. We now construct such a system for the QAP relation.

\todo{Explain why this NILP is non-interactive despite stemming from a Linear Interactive Proof (LIP): the CRS plays the role of the verifier's challenges, collapsing the interaction into a single message. Discuss the connection between 2-move LIPs and NILPs more carefully.}

\subsubsection*{Setup}

Pick a random trapdoor $\tau = (\alpha, \beta, \gamma, \delta, x) \in \F_p^{*5}$ and define the common reference string $\sigma$ as the collection of field evaluations
\[
\sigma =
\left(
\begin{aligned}
&\alpha,\; \beta,\; \gamma,\; \delta,\;\{x^i\}_{i=0}^{n-1},\\[2pt]
&\left\{\frac{\beta u_i(x)+\alpha v_i(x)+w_i(x)}{\gamma}\right\}_{i=0}^{\ell},\\[4pt]
&\left\{\frac{\beta u_i(x)+\alpha v_i(x)+w_i(x)}{\delta}\right\}_{i=\ell+1}^{m},\\[4pt]
&\left\{\frac{x^i\, t(x)}{\delta}\right\}_{i=0}^{n-2}
\end{aligned}
\right).
\]
The first block enables evaluation of arbitrary polynomials at $x$. The second block encodes the public witness terms, scaled by $\gamma$. The third encodes the private witness terms, scaled by $\delta$. The last encodes powers of $x$ weighted by $t(x)/\delta$, allowing the prover to compute $h(x)\,t(x)/\delta$ as a linear combination.

\todo{Discuss what $\alpha$, $\beta$, $\gamma$, and $\delta$ each do exactly: why they appear in the CRS, how they enforce consistency between $A$, $B$, $C$, and why the $\gamma$/$\delta$ separation is necessary.}

\subsubsection*{Prover}

Given a satisfying assignment $(a_1, \ldots, a_m)$, the prover computes $h(X)$ from the QAP relation, then picks two random blinding scalars $r, s \leftarrow \F_p$ and sets
\begin{align*}
  A &\deq \alpha + U(x) + r\delta, \\
  B &\deq \beta  + V(x) + s\delta, \\
  C &\deq \frac{\displaystyle\sum_{i=\ell+1}^{m} a_i\bigl(\beta u_i(x)+\alpha v_i(x)+w_i(x)\bigr) + h(x)\,t(x)}{\delta} + As + rB - rs\delta.
\end{align*}
The proof is $\pi = (A, B, C)$.

\subsubsection*{Verifier}

The verifier checks the single equation
\[
  A \cdot B
  \;\stackrel{?}{=}\;
  \underbrace{\alpha\beta}_{\text{(i)}}
  \;+\;
  \underbrace{\sum_{i=0}^{\ell} a_i \cdot \frac{\beta u_i(x)+\alpha v_i(x)+w_i(x)}{\gamma}}_{\text{(ii)}} \cdot\, \gamma
  \;+\;
  \underbrace{C \cdot \delta}_{\text{(iii)}}.
\]
The three terms on the right play distinct roles.
\begin{itemize}[noitemsep, leftmargin=2em]
  \item[(i)] $\alpha\beta$ is a fixed anchor from the setup. Its presence forces $A$ and $B$ to contain genuine $\alpha$ and $\beta$ offsets respectively, preventing a prover from constructing $A$ and $B$ independently of the structure of $\sigma$.
  \item[(ii)] The sum over $i = 0, \ldots, \ell$ incorporates the \emph{public inputs} $a_0, \ldots, a_\ell$ --- the statement $\phi$ that the verifier knows. These are precomputed in the CRS scaled by $\gamma$ and $\gamma$ cancels out, leaving the verifier to evaluate this term using only public data. Note that $\alpha$ and $\beta$ also appear here, but they are \emph{setup parameters} embedded in the precomputed CRS entries, not the circuit's inputs.
  \item[(iii)] $C \cdot \delta$ is where the private witness $a_{\ell+1}, \ldots, a_m$ is encoded. The $\delta$ scaling keeps it algebraically separated from term (ii), preventing the prover from mixing public and private contributions.
\end{itemize}

\subsubsection*{Simulator}

Given the trapdoor $\tau$, the simulator ignores the witness entirely. It picks $A, B \leftarrow \F_p$ uniformly at random and computes
\[
  C \deq \frac{A B \;-\; \alpha\beta \;-\; \displaystyle\sum_{i=0}^{\ell} a_i\bigl(\beta u_i(x)+\alpha v_i(x)+w_i(x)\bigr)}{\delta}.
\]

\subsubsection*{Why this achieves zero knowledge}

The blinding scalars $r$ and $s$ make $A$ and $B$ uniformly distributed in $\F_p$, independent of the witness. Given uniform $A$ and $B$, the verification equation uniquely determines $C$. The real distribution of $(A, B, C)$ is therefore identical to the simulated one: two uniform field elements followed by their forced companion. This is perfect zero knowledge.

\subsubsection*{Why this is sound against affine strategies}

An affine prover strategy computes $(A, B, C)$ as linear combinations of $\sigma$. Viewing the verification equation as a polynomial identity in the formal indeterminates $\alpha, \beta, \gamma, \delta, x$ and applying Schwartz--Zippel, any affine strategy that passes verification with non-negligible probability must satisfy the identity formally. A careful analysis of the monomials (particularly those involving $\alpha\beta$, $1/\gamma^2$, and $1/\delta^2$) forces the coefficients to encode a consistent assignment, from which the witness $(a_{\ell+1}, \ldots, a_m)$ can be extracted.

\todo{Add the full monomial analysis proof here, following Theorem~1 of the Groth (2016) paper.}

%% -----------------------------------------------------------
\subsection{Compiling to Pairings: The Generic Group Model}
%% -----------------------------------------------------------

The NILP works over a field, but we need a cryptographic argument --- the prover should not be able to read off $x$ from $\sigma$. The compilation step replaces every field element in $\sigma$ with its bracket encoding, and $(A, B, C)$ become $(\gone{A}, \gtwo{B}, \gone{C})$.

The compiled verification equation is:
\[
  \gone{A} \cdot \gtwo{B}
  \;=\;
  \gone{\alpha}\cdot\gtwo{\beta}
  \;+\;
  \sum_{i=0}^{\ell} a_i\,
    \left[\frac{\beta u_i(x)+\alpha v_i(x)+w_i(x)}{\gamma}\right]_{\!1}
    \cdot\, \gtwo{\gamma}
  \;+\;
  \gone{C}\cdot\gtwo{\delta}.
\]
The prover computes $\gone{A}$, $\gone{C}$, $\gtwo{B}$ as linear combinations of the CRS elements --- this is exactly the NILP prover running in the exponent. The verifier checks one pairing product equation using three pairings in total.

\textbf{Lifting soundness via the generic group model.} In the generic group model, a prover can only use the group elements in $\sigma$ through addition (and scalar multiplication) --- it cannot read their discrete logarithms. Any computation it performs on $\gone{\sigma_1}$ and $\gtwo{\sigma_2}$ is therefore a sequence of affine operations, corresponding exactly to choosing a proof matrix $\Pi$ and computing $\Pi\sigma$. This is precisely an affine strategy in the NILP sense. Combined with a \emph{disclosure-freeness} property of the CRS --- which guarantees that the structure of $\sigma$ cannot be exploited to choose a non-affine strategy --- the NILP's soundness against affine strategies lifts to soundness against all generic adversaries.

%% -----------------------------------------------------------
\subsection*{Summary: The Groth16 Proof System}
%% -----------------------------------------------------------

\begin{tcolorbox}[enhanced, colback=blue!3!white, colframe=blue!45!black,
                  arc=2.5pt, boxrule=1pt, left=6pt, right=6pt, top=5pt, bottom=5pt]
\small
\textbf{Setup.} Sample $\tau = (\alpha,\beta,\gamma,\delta,x) \leftarrow \F_p^{*5}$.
Publish $\sigma = \bigl([\sigma_1]_1,\, [\sigma_2]_2\bigr)$ as described above. Keep $\tau$ secret.

\medskip
\textbf{Prove.} Given $(a_1,\ldots,a_m)$ satisfying the circuit, pick $r,s \leftarrow \F_p$ and output
$\pi = \bigl(\gone{A},\, \gone{C},\, \gtwo{B}\bigr) \in \G_1^2 \times \G_2$.

\medskip
\textbf{Verify.} Check:
$\gone{A} \cdot \gtwo{B} = \gone{\alpha}\cdot\gtwo{\beta} + \displaystyle\sum_{i=0}^{\ell} a_i\,\gone{(\cdot)_i}\cdot\gtwo{\gamma} + \gone{C}\cdot\gtwo{\delta}$.

\medskip
\textbf{Proof size:} \quad $2$ elements in $\G_1$, \quad $1$ element in $\G_2$ \quad $= 3$ group elements.\\
\textbf{Verification:} \quad $1$ pairing product equation requiring $3$ online pairings:
$e(\gone{A}, \gtwo{B})$,\; $e\!\left(\textstyle\sum a_i K_i,\, \gtwo{\gamma}\right)$,\; $e(\gone{C}, \gtwo{\delta})$.
The fourth term $\gT{\alpha\beta}$ is precomputed and stored in the verifier key, so no extra pairing is needed for it.
\end{tcolorbox}
 inside zero_knowledge.tex

%% ============================================================
\section{Groth16: A Pairing-Based SNARK}
%% ============================================================

Groth16 is a preprocessing SNARK for arithmetic circuit satisfiability. Its proof is only three group elements, verification is a single pairing equation, and it achieves perfect completeness and perfect zero knowledge. Soundness holds against adversaries that interact with the system via generic group operations.

We build the system in three layers: encode computation as a polynomial divisibility problem (Quadratic Arithmetic Program, QAP), construct an information-theoretic proof for it (Non-Interactive Linear Proof, NILP), then compile the NILP into group elements using pairings.

%% -----------------------------------------------------------
\subsection{Notation for Pairing Groups}
%% -----------------------------------------------------------

Let $(p, \G_1, \G_2, \GT, e, g, h)$ be an asymmetric bilinear group of prime order $p$. The generator of $\G_1$ is $g$, the generator of $\G_2$ is $h$, and $e : \G_1 \times \G_2 \to \GT$ is a bilinear pairing satisfying $e(g^a, h^b) = e(g,h)^{ab}$.

We represent group elements by their discrete logarithms using the bracket notation
\[
  \gone{a} \deq g^a \in \G_1, \qquad
  \gtwo{a} \deq h^a \in \G_2, \qquad
  \gT{a}   \deq e(g,h)^a \in \GT, \qquad a \in \Z_p.
\]
This notation makes linear algebra readable. For a known scalar $c$ or known matrix $M$:
\[
  \gone{a} + \gone{b} = \gone{a+b}, \qquad
  c \cdot \gone{a} = \gone{ca}, \qquad
  M \cdot \gone{\mathbf{a}} = \gone{M\mathbf{a}}.
\]
The pairing provides one level of multiplication across the two source groups:
\[
  \gone{a} \cdot \gtwo{b} \deq e\!\left(g^a, h^b\right) = \gT{ab}.
\]


%% -----------------------------------------------------------
\subsection{From Arithmetic Circuits to Quadratic Arithmetic Programs}
%% -----------------------------------------------------------

\subsubsection*{Arithmetic circuits and R1CS}

An arithmetic circuit over a finite field $\F_p$ consists of addition and multiplication gates over field elements. Any such circuit can be \emph{flattened}: every gate is rewritten as a single multiplicative constraint
\[
  \bigl(\textstyle\sum_i A_{q,i}\, a_i\bigr)
  \cdot
  \bigl(\textstyle\sum_i B_{q,i}\, a_i\bigr)
  =
  \textstyle\sum_i C_{q,i}\, a_i
  \qquad (q = 1,\ldots, n),
\]
where $z = (a_0, a_1, \ldots, a_m) \in \F_p^{m+1}$ is the \emph{assignment vector}: $a_0 = 1$ is a constant, $\phi = (a_1, \ldots, a_\ell)$ is the public statement, and $w = (a_{\ell+1}, \ldots, a_m)$ is the secret witness. The $n$ constraints are collected into three matrices $A, B, C \in \F_p^{n \times (m+1)}$ --- this is the \emph{rank-1 constraint system} (R1CS). Arithmetic circuit satisfiability is NP-complete: any NP statement can be compiled into an R1CS instance over a suitable field, so a proof system for R1CS is a proof system for all of NP.

\subsubsection*{Quadratic arithmetic programs}

Checking $n$ separate scalar equations one by one would be costly. The key algebraic trick is to compress them into a single polynomial divisibility test.

Fix $n$ distinct field elements $r_1, \ldots, r_n \in \F_p$, one for each constraint. Define the vanishing polynomial
\[
  t(X) = \prod_{q=1}^{n}(X - r_q).
\]
For each variable index $i \in \{0, \ldots, m\}$, interpolate three polynomials $u_i, v_i, w_i \in \F_p[X]$ of degree at most $n-1$ such that
\[
  u_i(r_q) = A_{q,i}, \qquad v_i(r_q) = B_{q,i}, \qquad w_i(r_q) = C_{q,i} \qquad \forall\, q.
\]
\todo{Add a visualisation showing how the columns of the $A$, $B$, $C$ matrices are read off and interpolated into the polynomials $u_i$, $v_i$, $w_i$.}
Given an assignment $z$, define the aggregated polynomials
\[
  U(X) = \sum_{i=0}^{m} a_i\, u_i(X), \qquad
  V(X) = \sum_{i=0}^{m} a_i\, v_i(X), \qquad
  W(X) = \sum_{i=0}^{m} a_i\, w_i(X).
\]
At each constraint point $r_q$, the R1CS equation $U(r_q)\cdot V(r_q) = W(r_q)$ must hold. Since $t(X)$ vanishes precisely at $r_1, \ldots, r_n$, this is equivalent to
\[
  \boxed{t(X) \;\Big|\; U(X)\cdot V(X) - W(X).}
\]
Equivalently, there exists a degree-$(n-2)$ quotient polynomial $h(X)$ such that
\[
  U(X)\cdot V(X) = W(X) + h(X)\,t(X).
\]
This is the \emph{quadratic arithmetic program} (QAP) relation for the circuit. The advantage over R1CS is that $n$ separate equations have been replaced by one polynomial identity checkable at a single random point --- exactly the Schwartz--Zippel trick.

%% -----------------------------------------------------------
\subsection{A Non-Interactive Linear Proof for QAP}
%% -----------------------------------------------------------

A \emph{non-interactive linear proof} (NILP) is an information-theoretic proof system where the prover's messages are linear functions of a common reference string $\sigma$. We now construct such a system for the QAP relation.

\todo{Explain why this NILP is non-interactive despite stemming from a Linear Interactive Proof (LIP): the CRS plays the role of the verifier's challenges, collapsing the interaction into a single message. Discuss the connection between 2-move LIPs and NILPs more carefully.}

\subsubsection*{Setup}

Pick a random trapdoor $\tau = (\alpha, \beta, \gamma, \delta, x) \in \F_p^{*5}$ and define the common reference string $\sigma$ as the collection of field evaluations
\[
\sigma =
\left(
\begin{aligned}
&\alpha,\; \beta,\; \gamma,\; \delta,\;\{x^i\}_{i=0}^{n-1},\\[2pt]
&\left\{\frac{\beta u_i(x)+\alpha v_i(x)+w_i(x)}{\gamma}\right\}_{i=0}^{\ell},\\[4pt]
&\left\{\frac{\beta u_i(x)+\alpha v_i(x)+w_i(x)}{\delta}\right\}_{i=\ell+1}^{m},\\[4pt]
&\left\{\frac{x^i\, t(x)}{\delta}\right\}_{i=0}^{n-2}
\end{aligned}
\right).
\]
The first block enables evaluation of arbitrary polynomials at $x$. The second block encodes the public witness terms, scaled by $\gamma$. The third encodes the private witness terms, scaled by $\delta$. The last encodes powers of $x$ weighted by $t(x)/\delta$, allowing the prover to compute $h(x)\,t(x)/\delta$ as a linear combination.

\todo{Discuss what $\alpha$, $\beta$, $\gamma$, and $\delta$ each do exactly: why they appear in the CRS, how they enforce consistency between $A$, $B$, $C$, and why the $\gamma$/$\delta$ separation is necessary.}

\subsubsection*{Prover}

Given a satisfying assignment $(a_1, \ldots, a_m)$, the prover computes $h(X)$ from the QAP relation, then picks two random blinding scalars $r, s \leftarrow \F_p$ and sets
\begin{align*}
  A &\deq \alpha + U(x) + r\delta, \\
  B &\deq \beta  + V(x) + s\delta, \\
  C &\deq \frac{\displaystyle\sum_{i=\ell+1}^{m} a_i\bigl(\beta u_i(x)+\alpha v_i(x)+w_i(x)\bigr) + h(x)\,t(x)}{\delta} + As + rB - rs\delta.
\end{align*}
The proof is $\pi = (A, B, C)$.

\subsubsection*{Verifier}

The verifier checks the single equation
\[
  A \cdot B
  \;\stackrel{?}{=}\;
  \underbrace{\alpha\beta}_{\text{(i)}}
  \;+\;
  \underbrace{\sum_{i=0}^{\ell} a_i \cdot \frac{\beta u_i(x)+\alpha v_i(x)+w_i(x)}{\gamma}}_{\text{(ii)}} \cdot\, \gamma
  \;+\;
  \underbrace{C \cdot \delta}_{\text{(iii)}}.
\]
The three terms on the right play distinct roles.
\begin{itemize}[noitemsep, leftmargin=2em]
  \item[(i)] $\alpha\beta$ is a fixed anchor from the setup. Its presence forces $A$ and $B$ to contain genuine $\alpha$ and $\beta$ offsets respectively, preventing a prover from constructing $A$ and $B$ independently of the structure of $\sigma$.
  \item[(ii)] The sum over $i = 0, \ldots, \ell$ incorporates the \emph{public inputs} $a_0, \ldots, a_\ell$ --- the statement $\phi$ that the verifier knows. These are precomputed in the CRS scaled by $\gamma$ and $\gamma$ cancels out, leaving the verifier to evaluate this term using only public data. Note that $\alpha$ and $\beta$ also appear here, but they are \emph{setup parameters} embedded in the precomputed CRS entries, not the circuit's inputs.
  \item[(iii)] $C \cdot \delta$ is where the private witness $a_{\ell+1}, \ldots, a_m$ is encoded. The $\delta$ scaling keeps it algebraically separated from term (ii), preventing the prover from mixing public and private contributions.
\end{itemize}

\subsubsection*{Simulator}

Given the trapdoor $\tau$, the simulator ignores the witness entirely. It picks $A, B \leftarrow \F_p$ uniformly at random and computes
\[
  C \deq \frac{A B \;-\; \alpha\beta \;-\; \displaystyle\sum_{i=0}^{\ell} a_i\bigl(\beta u_i(x)+\alpha v_i(x)+w_i(x)\bigr)}{\delta}.
\]

\subsubsection*{Why this achieves zero knowledge}

The blinding scalars $r$ and $s$ make $A$ and $B$ uniformly distributed in $\F_p$, independent of the witness. Given uniform $A$ and $B$, the verification equation uniquely determines $C$. The real distribution of $(A, B, C)$ is therefore identical to the simulated one: two uniform field elements followed by their forced companion. This is perfect zero knowledge.

\subsubsection*{Why this is sound against affine strategies}

An affine prover strategy computes $(A, B, C)$ as linear combinations of $\sigma$. Viewing the verification equation as a polynomial identity in the formal indeterminates $\alpha, \beta, \gamma, \delta, x$ and applying Schwartz--Zippel, any affine strategy that passes verification with non-negligible probability must satisfy the identity formally. A careful analysis of the monomials (particularly those involving $\alpha\beta$, $1/\gamma^2$, and $1/\delta^2$) forces the coefficients to encode a consistent assignment, from which the witness $(a_{\ell+1}, \ldots, a_m)$ can be extracted.

\todo{Add the full monomial analysis proof here, following Theorem~1 of the Groth (2016) paper.}

%% -----------------------------------------------------------
\subsection{Compiling to Pairings: The Generic Group Model}
%% -----------------------------------------------------------

The NILP works over a field, but we need a cryptographic argument --- the prover should not be able to read off $x$ from $\sigma$. The compilation step replaces every field element in $\sigma$ with its bracket encoding, and $(A, B, C)$ become $(\gone{A}, \gtwo{B}, \gone{C})$.

The compiled verification equation is:
\[
  \gone{A} \cdot \gtwo{B}
  \;=\;
  \gone{\alpha}\cdot\gtwo{\beta}
  \;+\;
  \sum_{i=0}^{\ell} a_i\,
    \left[\frac{\beta u_i(x)+\alpha v_i(x)+w_i(x)}{\gamma}\right]_{\!1}
    \cdot\, \gtwo{\gamma}
  \;+\;
  \gone{C}\cdot\gtwo{\delta}.
\]
The prover computes $\gone{A}$, $\gone{C}$, $\gtwo{B}$ as linear combinations of the CRS elements --- this is exactly the NILP prover running in the exponent. The verifier checks one pairing product equation using three pairings in total.

\textbf{Lifting soundness via the generic group model.} In the generic group model, a prover can only use the group elements in $\sigma$ through addition (and scalar multiplication) --- it cannot read their discrete logarithms. Any computation it performs on $\gone{\sigma_1}$ and $\gtwo{\sigma_2}$ is therefore a sequence of affine operations, corresponding exactly to choosing a proof matrix $\Pi$ and computing $\Pi\sigma$. This is precisely an affine strategy in the NILP sense. Combined with a \emph{disclosure-freeness} property of the CRS --- which guarantees that the structure of $\sigma$ cannot be exploited to choose a non-affine strategy --- the NILP's soundness against affine strategies lifts to soundness against all generic adversaries.

%% -----------------------------------------------------------
\subsection*{Summary: The Groth16 Proof System}
%% -----------------------------------------------------------

\begin{tcolorbox}[enhanced, colback=blue!3!white, colframe=blue!45!black,
                  arc=2.5pt, boxrule=1pt, left=6pt, right=6pt, top=5pt, bottom=5pt]
\small
\textbf{Setup.} Sample $\tau = (\alpha,\beta,\gamma,\delta,x) \leftarrow \F_p^{*5}$.
Publish $\sigma = \bigl([\sigma_1]_1,\, [\sigma_2]_2\bigr)$ as described above. Keep $\tau$ secret.

\medskip
\textbf{Prove.} Given $(a_1,\ldots,a_m)$ satisfying the circuit, pick $r,s \leftarrow \F_p$ and output
$\pi = \bigl(\gone{A},\, \gone{C},\, \gtwo{B}\bigr) \in \G_1^2 \times \G_2$.

\medskip
\textbf{Verify.} Check:
$\gone{A} \cdot \gtwo{B} = \gone{\alpha}\cdot\gtwo{\beta} + \displaystyle\sum_{i=0}^{\ell} a_i\,\gone{(\cdot)_i}\cdot\gtwo{\gamma} + \gone{C}\cdot\gtwo{\delta}$.

\medskip
\textbf{Proof size:} \quad $2$ elements in $\G_1$, \quad $1$ element in $\G_2$ \quad $= 3$ group elements.\\
\textbf{Verification:} \quad $1$ pairing product equation requiring $3$ online pairings:
$e(\gone{A}, \gtwo{B})$,\; $e\!\left(\textstyle\sum a_i K_i,\, \gtwo{\gamma}\right)$,\; $e(\gone{C}, \gtwo{\delta})$.
The fourth term $\gT{\alpha\beta}$ is precomputed and stored in the verifier key, so no extra pairing is needed for it.
\end{tcolorbox}
